% Ch. 27 : programmation classique contre apprentissage automatique
\documentclass[border=6pt]{standalone}
\usepackage{coursfig}
\begin{document}
\begin{tikzpicture}[x=1mm,y=1mm]
  % programmation classique
  \node[box=Enc, text width=30mm] (r1) at (0,8) {\textbf{règles}\sub{écrites par vous}};
  \node[box=Rule, fill=white, text width=30mm] (d1) at (0,-8) {\textbf{données}};
  \node[keybox=Enc, text width=30mm, minimum height=14mm] (p1) at (52,0) {programme};
  \node[box=Sea, text width=30mm] (o1) at (104,0) {\textbf{réponses}};
  \draw[flow] (r1.east) -- (p1.west); \draw[flow] (d1.east) -- (p1.west); \draw[flow] (p1) -- (o1);
  \node[ftitle, anchor=south west] at (-16,20) {PROGRAMMATION CLASSIQUE};
  % apprentissage
  \begin{scope}[yshift=-46mm]
    \node[box=Rule, fill=white, text width=30mm] (d2) at (0,8) {\textbf{données}};
    \node[box=Sea, text width=30mm] (a2) at (0,-8) {\textbf{réponses}\sub{attendues (étiquettes)}};
    \node[keybox=Ocre, text width=30mm, minimum height=14mm] (p2) at (52,0) {apprentissage};
    \node[box=Enc, text width=30mm] (m2) at (104,0) {\textbf{modèle}\sub{des règles apprises}};
    \draw[flow] (d2.east) -- (p2.west); \draw[flow] (a2.east) -- (p2.west); \draw[flow] (p2) -- (m2);
    \node[ftitle, anchor=south west] at (-16,20) {APPRENTISSAGE AUTOMATIQUE};
  \end{scope}
  \node[note, anchor=west, text width=46mm] at (124,-46) {le comportement du système dépend désormais des \textbf{données}, pas seulement du code};
\end{tikzpicture}
\end{document}
